%!TEX root = thesis.tex

Baseline drift in chromatographic signals corrupts peak area integration, the standard method for analyte quantification, yet existing correction methods require per-signal parameter tuning (classical optimization) or sacrifice interpretability (deep learning). This thesis proposes LBEADS-NET, a neural network that bridges this gap through \emph{algorithm unrolling}: each iteration of the classical BEADS (Baseline Estimation And Denoising using Sparsity) algorithm is mapped to a differentiable network layer with learnable parameters, creating a 20-layer architecture with 100 learned scalar parameters that inherits the optimization structure of BEADS while gaining adaptivity through end-to-end training. The ISTA-style proximal gradient formulation, adopted after a conjugate gradient variant exhibited gradient vanishing at effective depth 96, enables stable training across all layers. A multi-stage curriculum manages an 11-term composite loss function designed to suppress \emph{baseline leakage}, a phenomenon in which the $\ell_1$ sparsity penalty causes peak energy to be absorbed into the baseline estimate in dense-peak regions. Progressive experiments show that all LBEADS-NET variants outperform classical BEADS with fixed parameters on peak reconstruction error (${\sim}2\times$ improvement) without per-signal tuning. However, the Baseline Leakage Index plateaus at approximately 0.59 across all loss configurations, revealing a fundamental limit: loss-based mitigations are regularizers on a sparsity-biased architecture, and no amount of loss engineering overcomes the leakage when the sparsity assumption is violated. This finding, corroborated by concurrent work on unrolled sparse-recovery networks, generalizes to any algorithm-unrolled model inheriting $\ell_1$ priors. The thesis further contributes analyses of the autograd boundary (non-learnable filter cutoff frequency $f_c$), training/inference length mismatch, and the over-regularization effect observed when 12 competing loss terms exceed the training budget. Learned parameters exhibit emergent layer specialization, increasing sparsity weight and non-monotonic asymmetry ratio across depth, demonstrating that algorithm unrolling discovers iteration schedules inaccessible to the fixed-parameter classical algorithm.
